\documentclass[11pt]{article}

\usepackage[margin=1in]{geometry}
\usepackage{amsmath, amssymb, amsfonts}
\usepackage{booktabs}
\usepackage{enumitem}
\usepackage{graphicx}
\usepackage{hyperref}
\usepackage{xcolor}
\graphicspath{{../../artifacts/}}

\title{Methodology Evaluation Framework for Classification Studies}
\author{kinematic-classifier-sandbox}
\date{\today}

\begin{document}

\maketitle

\section{Purpose}

This document complements the posterior update note. The posterior note explains
how class evidence is accumulated. This document explains how the repository
evaluates whether that machinery is trustworthy, informative, and methodologically
useful.

The main questions are:
\begin{itemize}[leftmargin=1.5em]
    \item How prior-sensitive are the decisions?
    \item How separable are the classes in the current feature space?
    \item How should confusion matrices be interpreted?
    \item How do we decide whether a corpus is too easy, too biased, or
    under-exercised?
\end{itemize}

Primary implementation surfaces covered explicitly by this note:
\begin{itemize}[leftmargin=1.5em]
    \item \texttt{prior\_sensitivity\_analysis.py}
    \item \texttt{feature\_analysis.py}
    \item \texttt{pca\_analysis.py}
    \item \texttt{short\_horizon\_identifiability.py}
    \item \texttt{corpus\_adequacy\_audit.py}
    \item \texttt{coverage\_report.py}
    \item \texttt{inspection\_bundle.py}
    \item \texttt{generic\_feature\_taxonomy.py}
    \item \texttt{class\_validity.py}
    \item \texttt{generated\_corpus\_features.py}
    \item \texttt{corpus\_classifier\_scoring.py}
\end{itemize}

\section{Problem And Notation}

The central evaluation question is:
\begin{equation}
    \text{Given a classifier, feature family, and corpus, what evidence shows whether apparent success or failure is real?}
\end{equation}

The most important symbols are:
\begin{itemize}[leftmargin=1.5em]
    \item \(p(A)\): prior probability of class \(A\),
    \item \(L_t(c)\): likelihood term for class \(c\) at time \(t\),
    \item \(C_{a,b}\): confusion count for true class \(a\) predicted as \(b\),
    \item \(\phi_t\): feature vector at time \(t\),
    \item \(\mathrm{AUC}(a,b)\): pairwise AUC for classes \(a\) and \(b\),
    \item \(O(a,b)\): overlap estimate for classes \(a\) and \(b\).
\end{itemize}

\section{Evaluation Layers}

The repo is increasingly using the following evaluation stack:
\begin{equation}
    \text{corpus}
    \rightarrow
    \text{features or filters}
    \rightarrow
    \text{evidence}
    \rightarrow
    \text{posterior}
    \rightarrow
    \text{metrics}
    \rightarrow
    \text{artifact bundles}.
\end{equation}

Different artifact families probe different layers:
\begin{itemize}[leftmargin=1.5em]
    \item prior sensitivity probes the posterior layer,
    \item AUC, overlap, and identifiability probe the feature-evidence layer,
    \item confusion matrices probe end-to-end decision behavior,
    \item corpus adequacy probes the data-generation layer.
\end{itemize}

That separation matters. A poor confusion matrix does not automatically imply a
bad posterior updater. It may indicate weak features or a biased corpus instead.

\section{Assumptions Behind The Evaluation Stack}

The evaluation layer assumes:
\begin{itemize}[leftmargin=1.5em]
    \item no single metric is sufficient for methodological judgment,
    \item priors, AUC, confusion, and adequacy measure different things,
    \item feature-space geometry should be inspected before blaming the posterior updater,
    \item corpus leakage can invalidate apparently good classifier results.
\end{itemize}

\section{How We Evaluate Priors}

\subsection{Prior Sweep}

For a binary decision between classes $A$ and $B$, the prior-sensitivity study
sweeps:
\begin{equation}
    p(A) \in [0.05, 0.95],
    \qquad
    p(B) = 1 - p(A).
\end{equation}

At each prior setting, the classifier is rerun and the final posterior and hard
decision are recorded.

\subsection{Flip Thresholds and Fragility}

The key outputs are:
\begin{itemize}[leftmargin=1.5em]
    \item the smallest prior shift needed to flip the final class,
    \item the smallest log-prior shift needed to flip the final class,
    \item the fraction of trajectories that flip under a moderate perturbation.
\end{itemize}

Interpretation is straightforward:
\begin{itemize}[leftmargin=1.5em]
    \item large required shifts imply evidence-dominated decisions,
    \item small required shifts imply prior-fragile decisions.
\end{itemize}

\subsection{Odds Decomposition}

For accumulator-style models, the useful identity is:
\begin{equation}
    \log \frac{p(A \mid z)}{p(B \mid z)}
    =
    \log \frac{p(A)}{p(B)}
    +
    \sum_t \log \frac{L_A(z_t)}{L_B(z_t)}.
\end{equation}

This makes the interpretation cleaner:
\begin{itemize}[leftmargin=1.5em]
    \item strong cumulative likelihood separation reduces prior influence,
    \item weak cumulative separation allows the prior to dominate.
\end{itemize}

\subsection{Current Repo Interpretation}

The current prior-sensitivity artifacts show the desired split:
\begin{itemize}[leftmargin=1.5em]
    \item easy trajectories are stable under moderate prior changes,
    \item ambiguous trajectories expose prior-dominant regions,
    \item the fragility view is more informative than a simple flip/not-flip
    label because it shows how close a trajectory is to changing class.
\end{itemize}

\begin{figure}[htbp]
    \centering
    \includegraphics[width=0.48\textwidth]{prior_sensitivity_v1/trajectory_prior_fragility_overview.png}
    \hfill
    \includegraphics[width=0.48\textwidth]{prior_sensitivity_v1/pairwise_flip_threshold_heatmap.png}
    \caption{Two complementary prior-sensitivity views: trajectory-level
    fragility and pairwise flip-threshold structure.}
\end{figure}

\section{How We Evaluate Separability}

\subsection{Pairwise AUC}

For a class pair $(a,b)$, pairwise AUC measures ranking separation:
\begin{equation}
    \mathrm{AUC}(a,b) \in [0.5, 1.0].
\end{equation}

Interpretation:
\begin{itemize}[leftmargin=1.5em]
    \item values near $1.0$ indicate strong ordering separation,
    \item values near $0.5$ indicate little or no ranking separation.
\end{itemize}

But AUC alone is not enough. AUC can be high because the corpus is unrealistically
easy.

\subsection{Overlap and Distance}

The repo also tracks overlap and pairwise distance summaries. Their methodological
roles are:
\begin{itemize}[leftmargin=1.5em]
    \item low overlap suggests distinct feature-space geometry,
    \item high overlap suggests intrinsic ambiguity unless richer features or
    longer histories are introduced,
    \item distance measures provide a geometry-aware complement to raw confusion.
\end{itemize}

\subsection{Oracle Versus Classifier Interpretation}

The important diagnostic split is:
\begin{itemize}[leftmargin=1.5em]
    \item low AUC and high overlap suggest a feature problem,
    \item good AUC but poor final confusion suggests a posterior, calibration,
    or model problem,
    \item perfect AUC on a supposedly hard pair suggests a corpus-design
    problem.
\end{itemize}

\subsection{Current Repo Interpretation}

The current abstract inspection bundle reports:
\begin{itemize}[leftmargin=1.5em]
    \item recommended feature set: \texttt{all\_engineered},
    \item hardest current class pair:
    \texttt{bounded\_acceleration vs maneuver}.
\end{itemize}

This is useful because the bundle combines:
\begin{itemize}[leftmargin=1.5em]
    \item average pairwise AUC,
    \item worst-case pairwise AUC,
    \item average overlap,
    \item worst-case overlap,
    \item and adequacy gates.
\end{itemize}

\begin{figure}[htbp]
    \centering
    \includegraphics[width=0.48\textwidth]{abstract_inspection_v1/feature_set_inspection_summary.png}
    \hfill
    \includegraphics[width=0.48\textwidth]{abstract_inspection_v1/hardest_class_pairs.png}
    \caption{Cross-feature-set inspection summary and hardest current class
    boundaries in the abstract inspection bundle.}
\end{figure}

\begin{figure}[htbp]
    \centering
    \includegraphics[width=0.48\textwidth]{feature_analysis_v1/class_confusability_heatmap.png}
    \hfill
    \includegraphics[width=0.48\textwidth]{feature_analysis_v1/pairwise_overlap_heatmap.png}
    \caption{Confusability and overlap views from the feature-space analysis
    layer.}
\end{figure}

\subsection{Generated Corpus Feature Integration}

\texttt{generated\_corpus\_features.py} is the bridge that takes
objective-driven, backend-generated trajectories and routes them through the
real feature extraction pipeline using class-validity-adjusted labels.

The methodological flow is
\begin{equation}
    \text{generated candidate}
    \rightarrow
    \text{executed trajectory}
    \rightarrow
    \text{validity-adjusted label}
    \rightarrow
    \text{feature extraction}
    \rightarrow
    \text{excitation and separability artifacts}.
\end{equation}

Methodologically, that matters because it closes the loop:
\begin{itemize}[leftmargin=1.5em]
    \item generated candidates are not only judged by proxy score rows,
    \item relabeled or ambiguous trajectories can still be inspected in feature
    space,
    \item corpus-generation work can be audited with the same feature and
    excitation machinery used elsewhere.
\end{itemize}

\subsection{Corpus-Conditioned Classifier Scoring}

\texttt{corpus\_classifier\_scoring.py} sits one layer downstream from
generated-corpus feature integration. It rebuilds pointwise, accumulator,
windowed, and Kalman-style scoring surfaces directly on generated-corpus
records so the repo can inspect disagreement, posterior entropy, confidence
margin, and prior sensitivity on the generated corpus itself.

Its stress proxy is approximately
\begin{equation}
    \text{stress}
    \approx
    0.5 \cdot (1 - \text{margin})
    + 0.5 \cdot \text{entropy}
    + 0.35 \cdot \mathbf{1}\{\text{final prediction wrong}\},
\end{equation}
which makes this layer a deliberate classifier-pressure surface rather than
just another export table.

\section{How We Interpret Confusion Matrices}

Confusion is an end-to-end decision artifact:
\begin{equation}
    C_{a,b} = \#\{\text{true class } a, \text{ predicted class } b\}.
\end{equation}

It is valuable because it shows what the decision system actually does. But it
is not a sufficient explanation by itself.

\subsection{Why Confusion Can Mislead}

The same off-diagonal confusion may come from:
\begin{itemize}[leftmargin=1.5em]
    \item feature insufficiency,
    \item poor likelihood shape,
    \item over-aggressive prior structure,
    \item weak within-class mode logic,
    \item or biased corpus construction.
\end{itemize}

That is why the repo pairs confusion with:
\begin{itemize}[leftmargin=1.5em]
    \item feature-vs-class matrices,
    \item entropy traces,
    \item pairwise AUC and overlap,
    \item prior-sensitivity studies.
\end{itemize}

\subsection{Multi-View Confusion}

The richer benchmark families already distinguish:
\begin{itemize}[leftmargin=1.5em]
    \item full-trajectory confusion,
    \item transient confusion,
    \item terminal confusion,
    \item phase confusion,
    \item feature-vs-class confusion.
\end{itemize}

This is methodologically stronger than a single final confusion matrix because
it localizes \emph{where} in time the error appears and whether it is tied to
missing feature anchors.

\subsection{Class Validity And Relabel Pressure}

\texttt{class\_validity.py} sits between corpus design and classifier
evaluation. It is not scoring posterior behavior directly. It is scoring whether
generated examples still resemble the class they were intended to instantiate.

The module computes telemetry-derived similarity scores for candidate classes
and then assigns label states such as:
\begin{itemize}[leftmargin=1.5em]
    \item \texttt{valid\_target\_class},
    \item \texttt{ambiguous},
    \item \texttt{invalid},
    \item \texttt{relabel\_candidate}.
\end{itemize}

Methodologically, this matters because some downstream confusion should be
blamed on weak or ambiguous class definitions rather than on the classifier
itself.

\section{How We Evaluate Corpus Quality}

\subsection{Balance Checks}

The corpus should be checked for:
\begin{itemize}[leftmargin=1.5em]
    \item class balance,
    \item scenario balance,
    \item duration balance,
    \item sample-count balance,
    \item noise and irregular-sampling coverage.
\end{itemize}

If those are class-linked, the classifier can appear stronger than it really is.

\subsection{Feature Excitation Coverage}

Each feature set should be excited across multiple tiers and classes. The
important question is not only whether a feature exists, but whether the corpus
actually exercises it often enough to judge its utility.

\subsection{Class-Pair Boundary Coverage}

Declared hard pairs should appear in the tiers where they are supposed to be
challenging:
\begin{itemize}[leftmargin=1.5em]
    \item boundary,
    \item adversarial,
    \item realistic,
    \item stress.
\end{itemize}

If a hard pair is missing or over-separated, then the study is not proving the
right thing.

\subsection{Leakage Audit}

The corpus should also be checked for covariate-only shortcuts such as:
\begin{itemize}[leftmargin=1.5em]
    \item duration,
    \item sample count,
    \item sampling irregularity,
    \item noise level,
    \item outlier fraction.
\end{itemize}

If a covariate alone predicts class membership too well, the classifier may be
benefiting from corpus bias instead of motion structure.

\subsection{Current Repo Interpretation}

The current corpus adequacy audit is already enforceable rather than merely
descriptive:
\begin{itemize}[leftmargin=1.5em]
    \item overall status is currently \texttt{fail},
    \item the main red finding is that
    \texttt{constant\_acceleration vs maneuver} is over-separated even though it
    was declared hard,
    \item two yellow findings indicate class-linked \texttt{std\_dt} and
    \texttt{sampling\_irregularity} imbalance.
\end{itemize}

\begin{figure}[htbp]
    \centering
    \includegraphics[width=0.48\textwidth]{corpus_adequacy_audit_v1/class_pair_coverage_heatmap.png}
    \hfill
    \includegraphics[width=0.48\textwidth]{corpus_adequacy_audit_v1/covariate_leakage_audit.png}
    \caption{Corpus quality views: class-pair boundary coverage and covariate
    leakage structure.}
\end{figure}

\section{How These Evaluations Connect}

The intended interpretation order is:
\begin{enumerate}[leftmargin=1.8em]
    \item check corpus adequacy first,
    \item then inspect feature excitation and pairwise separability,
    \item then inspect posterior behavior and prior sensitivity,
    \item then interpret confusion matrices.
\end{enumerate}

That ordering matters because otherwise:
\begin{itemize}[leftmargin=1.5em]
    \item a good confusion matrix can hide a bad corpus,
    \item and a bad confusion matrix can be misdiagnosed as a classifier problem
    when it is really a feature or data problem.
\end{itemize}

\section{Current Working Standard for ``Good Enough''}

Within this repo, a study component is behaving well when:
\begin{itemize}[leftmargin=1.5em]
    \item priors do not dominate easy decisions,
    \item hard class pairs are actually present and remain challenging,
    \item feature sets can be compared with stable AUC and overlap summaries,
    \item confusion can be explained by supporting diagnostics,
    \item the corpus does not leak labels through easy covariates.
\end{itemize}

That is a stronger standard than simply asking whether one top-line accuracy
number is high.

\section{Artifact Map}

The main artifact families for this evaluation framework are:
\begin{itemize}[leftmargin=1.5em]
    \item priors:
    \texttt{artifacts/prior\_sensitivity\_v1/}
    \item feature-space separability:
    \texttt{artifacts/feature\_analysis\_v1/} and
    \texttt{artifacts/abstract\_inspection\_v1/}
    \item corpus quality:
    \texttt{artifacts/corpus\_adequacy\_audit\_v1/} and
    \texttt{artifacts/coverage\_report\_v1/}
\end{itemize}

This document should be read alongside the posterior note:
\begin{itemize}[leftmargin=1.5em]
    \item \texttt{docs/surveys/posterior\_update\_math.tex}
\end{itemize}

That keeps posterior mechanics separate from evaluation methodology.

\end{document}
